\begin{abstract}
Conventional volumetric fluorescence imaging acquires images at successive focal planes.
Depth-dependent illumination offers an alternative: encode the volume into planar
measurements and reconstruct it at a fixed detection focus. We combine angle-scanned
phase illumination with multi-slice plug-and-play reconstruction. Varying the input angle
at a fixed phase encoder changes the volumetric excitation while the specimen and
detection focus remain stationary. Reconstruction couples measurement consistency with
a pretrained two-dimensional denoiser applied in three orthogonal slice orientations,
recovering a shared volume without task-specific three-dimensional network training.
Numerical evaluation with reference volumes of mouse pyramidal neurons with dendritic
arbors demonstrates recovery over a
\(30.72\times30.72\times10.24\,\mu\mathrm m^3\) volume with isotropic
\(0.64\,\mu\mathrm m\) reconstruction sampling. Dominant dendritic structures and their
depth ordering remain recognisable in representative reconstructions at eight-fold
measurement compression. The framework supports compressive volumetric fluorescence
recovery without axial specimen scanning or detection refocusing.
\end{abstract}
